\documentclass[10pt,a4paper]{amsart}
\usepackage[margin=19mm]{geometry}
\usepackage{amsmath,amssymb,mathtools}
\usepackage[scr=rsfs,cal=euler]{mathalfa}
\usepackage{xfrac}
\usepackage[unicode,hidelinks,bookmarks=false]{hyperref}
\newcommand{\ie}{i.e.\ }
\newcommand{\cf}{cf.\ }
\newcommand{\resp}{respectively\ }
\newcommand{\Subsection}{\subsection}
\providecommand{\nobreakdash}{\nobreak\hbox{-}}
\setlength{\emergencystretch}{2em}
\multlinegap0pt
\usepackage{tikz}
\usetikzlibrary{cd}
\usepackage[shortlabels]{enumitem}
\usepackage{mathtools}
\usepackage{amssymb}
\usepackage{xcolor}
\newtheorem{lemma}{Lemma}
\newtheorem{definition}{Definition}
\newtheorem{proposition}{Proposition}
\usepackage{cleveref}
\crefname{lemma}{Lemma}{Lemmas}
\crefname{definition}{Definition}{Definitions}
\crefname{proposition}{Proposition}{Propositions}
\newcommand{\Adm}{\operatorname{Adm}}
\newcommand{\Coh}{\operatorname{Coh}}
\newcommand{\Isoc}{\operatorname{Isoc}}
\newcommand{\Lie}{\operatorname{Lie}}
\newcommand{\fstr}{\mathrm{fstr}}
\title{Arithmetic Kashiwara Regularity and Orbit Classification for Filtered Strongly Equivariant $\mathcal{D}^{\dagger}$-Modules}
\author{Andr\'es Sarrazola-Alzate}
\date{Source: arXiv:2606.05547v1 (CC BY 4.0)}
\begin{document}
\maketitle

\noindent\emph{Source and licence.} Arithmetic Kashiwara Regularity and Orbit Classification for Filtered Strongly Equivariant $\mathcal{D}^{\dagger}$-Modules. Version arXiv:2606.05547v1 (CC BY 4.0), distributed by arXiv under the Creative Commons Attribution 4.0 International licence, http://creativecommons.org/licenses/by/4.0/, as recorded in the arXiv licence metadata for this exact version and shown on the abstract page https://arxiv.org/abs/2606.05547v1. Full licence notice: \texttt{licenses/arxiv-2606.05547-CC-BY-4.0.txt}.\par\smallskip

\noindent\textit{Editorial scope.} Counted unit: the article's proposition prop:eq-intermediate-extension (source0.tex lines 3098--3114). The article's complete original proof of it is delivered with it (source0.tex lines 3116--3225, one \texttt{proof} environment of 110 lines). No mathematics is written, repaired or re-derived: the statement and the proof are the article's own lines, in the article's own notation, and no retained line is substituted or re-flowed.  Retained uncounted as the source-local context the proof needs: lines 1632--1649; lines 1754--1800; lines 2958--2980; lines 3019--3025; lines 3035--3042; lines 3054--3074.  Nothing was omitted from any retained span, so the package carries no omission record.  No retained line contains a citation, so the wrapper carries no bibliography and no bibliography entry.  The retained lines contain the article's own diagram code, which the wrapper typesets with the standard tikz packages; no image file is delivered and the package declares no asset dependency.  Editorial formatting only, in addition: an AMS \texttt{amsart} preamble that restates the article's own declarations and macros used by the retained lines, namely the environments \texttt{lemma}, \texttt{definition}, \texttt{proposition} on separate counters, together with the standard packages those lines need.  The retained environments print in this document as Lemma 1, Definition 1, Proposition 1 in order of appearance, whereas the article prints the same items with the numbering of its own full document.  The complete native source of the article as distributed in the arXiv e-print for this exact version is bundled unchanged as \texttt{sources/202124/source0.tex}.
\begin{lemma}[Differentiation of equivariant structures]
\label{lem:differentiation-equivariant-structures}
Let \((\mathcal M,\alpha_{\mathcal M})\) be weakly
\(\mathfrak K\)-equivariant.  The isomorphism
\[
  \alpha_{\mathcal M}:a^\sharp\mathcal M\xrightarrow{\sim}p^\sharp\mathcal M
\]
induces a \(K\)-linear Lie algebra action
\[
  d\alpha_{\mathcal M}:\mathfrak k_K\longrightarrow
  \operatorname{End}_{\mathbb Q}(\mathcal M),
  \qquad
  \xi\longmapsto (m\mapsto \xi\cdot_\alpha m),
\]
which is functorial in weakly equivariant morphisms.  Moreover, kernels,
cokernels, images and coimages of weakly equivariant morphisms inherit the
induced infinitesimal action by restriction or passage to quotient.
\end{lemma}

\begin{definition}[Filtered strong equivariance]
\label{def:strong-equivariance}
\textup{\textbf{Filtered strong equivariance}} is the following finite-level
condition.  A weakly \(\mathfrak K\)-equivariant coherent
\(\mathscr D^\dagger_{\mathfrak X,\mathbb Q}\)-module
\((\mathcal M,\alpha_{\mathcal M})\) is called filtered strongly
\(\mathfrak K\)-equivariant if it satisfies \textup{(SE1)} and
\textup{(SE2)} below.
\begin{enumerate}[label=\textup{(SE\arabic*)}]
\item \emph{Infinitesimal compatibility.}  For every
\(\xi\in\mathfrak k_K\), the infinitesimal action induced by
\(\alpha_{\mathcal M}\) agrees with the action of the corresponding
fundamental vector field:
\[
  \xi\cdot_{\alpha}m=\mu_{\mathfrak X}(\xi)m.
\]
\item \emph{Filtered finite-level realizability.} There is a finite
affine open covering \(\{\mathfrak U_i\}_{i\in I}\) of
\(\mathfrak X\), an integer \(m_0\), and a finite \(K\)-basis
\(\xi_1,\ldots,\xi_r\) of the finite-dimensional Lie algebra
\(\mathfrak k_K=\Lie(\mathcal K)\otimes_{\mathcal V}K\) such that, for every
\(i\), the restriction \(\mathcal M|_{\mathfrak U_i}\) admits a coherent
level-\(m_0\) model \(\mathcal M_i^{(m_0)}\) with a good filtration
\(F_\bullet\mathcal M_i^{(m_0)}\), and the differentiated operators
\(d\alpha(\xi_a)\) are realized on this model and satisfy
\[
  d\alpha(\xi_a)\bigl(F_n\mathcal M_i^{(m_0)}\bigr)
  \subseteq F_n\mathcal M_i^{(m_0)}
  \qquad(a=1,\ldots,r,\ n\in\mathbb Z).
\]
Here ``realized on this model'' means that the endomorphism of the dagger
module induced by \(d\alpha(\xi_a)\) restricts to an endomorphism of the
chosen coherent
\(\widehat{\mathscr D}^{(m_0)}_{\mathfrak U_i,\mathbb Q}\)-model.  By
\(K\)-linearity the same preservation property then holds for every
\(\xi\in\mathfrak k_K\).  The level \(m_0\) is uniform on the finite cover,
which is the minimal form needed for the global characteristic-variety
argument.  We do not require compatible models at every larger level; if a
larger level is needed, one may extend scalars from the chosen level model, but
this extra uniformity is not part of the definition.
\end{enumerate}
We denote this exact category by
\[
  \Coh^{\fstr}_{\mathfrak K}
  \bigl(\mathscr D^\dagger_{\mathfrak X,\mathbb Q}\bigr).
\]
\end{definition}

\begin{lemma}[Smooth pullback and intermediate extension]
\label{lem:smooth-pullback-intermediate-extension}
Let \(f:Y'\to Y\) be a smooth morphism of smooth varieties over \(k\), and let
\(j:U\hookrightarrow Y\) be a locally closed immersion. Form the cartesian
square
\[
\begin{tikzcd}
  U' \arrow[r,"j'"] \arrow[d,"f_U"'] &
  Y' \arrow[d,"f"] \\
  U \arrow[r,"j"] &
  Y.
\end{tikzcd}
\]
Let \(E\) be a coefficient object on \((U,\overline U)\) belonging to the
overholonomic heart; in particular, this applies to an overconvergent
\(F\)-isocrystal. Then there is a canonical isomorphism
\[
  f^\sharp j_{!+}(E)
  \simeq
  j'_{!+}\bigl(f_U^*E\bigr),
\]
functorial in \(E\).
\end{lemma}

\begin{lemma}[Boundary vanishing for endomorphism differences]
\label{lem:boundary-vanishing-endomorphism-difference}
Let \(N=j_{O,!+}(E)\), with \(E\) in the overholonomic heart on
\((O,\overline O)\). Suppose that \(\delta:N\to N\) is a
\(\mathscr D^\dagger_{\mathfrak X,\mathbb Q}\)-linear endomorphism whose
restriction to \(O\) is zero. Then \(\delta=0\).
\end{lemma}

\begin{lemma}[Boundary vanishing after smooth pullback]
\label{lem:boundary-vanishing-smooth-pullback}
Let \(N=j_{O,!+}(E)\), with \(E\) in the overholonomic heart on
\((O,\overline O)\). Let \(f:Y'\to X_s\) be a smooth morphism and put
\(O':=f^{-1}(O)\). Then any
\(\mathscr D^\dagger\)-linear endomorphism of \(f^\sharp N\) whose restriction
to \(O'\) is zero is itself zero.
\end{lemma}

\begin{lemma}[Infinitesimal functoriality of Berthelot operators]
\label{lem:infinitesimal-functoriality-operators}
Let \(\mathfrak K\) act on \(\mathfrak X\), and let
\(\xi\in\mathfrak k_K\) with fundamental vector field \(v_\xi\). The induced
action of \(\mathfrak K\) on the finite-level sheaves
\(\widehat{\mathscr D}^{(m)}_{\mathfrak X,\mathbb Q}\) differentiates to
derivations
\[
  \xi_{\mathscr D}^{(m)}:
  \widehat{\mathscr D}^{(m)}_{\mathfrak X,\mathbb Q}
  \longrightarrow
  \widehat{\mathscr D}^{(m)}_{\mathfrak X,\mathbb Q}
\]
compatible with the transition maps in \(m\). For every local section
\(P\) of \(\widehat{\mathscr D}^{(m)}_{\mathfrak X,\mathbb Q}\), one has
\[
  [v_\xi,P]=\xi_{\mathscr D}^{(m)}(P).
\]
Passing to the inductive limit gives the corresponding identity on
\(\mathscr D^\dagger_{\mathfrak X,\mathbb Q}\).
\end{lemma}

\begin{proposition}[Equivariance of intermediate extension]
\label{prop:eq-intermediate-extension}
Let
\[
  E\in F\text{-}\Isoc^{\dagger}_{\mathcal K}(O,\overline O).
\]
Then \(j_{O,!+}(E)\) carries a natural Frobenius weak \(\mathfrak K\)-equivariant
structure satisfying the ordinary infinitesimal strong compatibility
\(d\alpha(\xi)=v_\xi\). If moreover
\([E]\in\Adm^{\Phi,\fstr}_{\mathcal K}(O,\overline O)\), then this middle
extension belongs to
\[
  F\text{-}\Coh^{\fstr}_{\mathfrak K}
  \bigl(\mathscr D^\dagger_{\mathfrak X,\mathbb Q}\bigr)
\]
by the defining filtered finite-level admissibility condition.
\end{proposition}

\begin{proof}
Let \(\beta_E:a_O^*E\xrightarrow{\sim}p_O^*E\) be the equivariant structure on
\(E\). Applying \Cref{lem:smooth-pullback-intermediate-extension} to the smooth
special-fiber action and projection maps gives canonical identifications
\[
  a^\sharp j_{O,!+}(E)
  \simeq
  j_{\mathcal K_s\times O,!+}(a_O^*E),
  \qquad
  p^\sharp j_{O,!+}(E)
  \simeq
  j_{\mathcal K_s\times O,!+}(p_O^*E).
\]
The isomorphism \(\beta_E\) therefore induces an isomorphism
\[
  \alpha_{j_{O,!+}(E)}:
  a^\sharp j_{O,!+}(E)
  \xrightarrow{\;\sim\;}
  p^\sharp j_{O,!+}(E),
\]
which is a weak \(\mathfrak K\)-equivariant structure. Frobenius compatibility
is inherited from the Frobenius compatibility of \(\beta_E\) and from the
functoriality of intermediate extension in the Frobenius overholonomic heart.

The identity and cocycle conditions are checked on the dense open orbit, where
they are precisely the identity and cocycle conditions for \(\beta_E\). For the
identity condition this gives two endomorphisms of \(j_{O,!+}(E)\) which
coincide after restriction to \(O\), and
\Cref{lem:boundary-vanishing-endomorphism-difference} forces them to coincide
on \(X_s\). For the cocycle condition the two compositions live after smooth
pullback to \(\mathcal K_s\times\mathcal K_s\times X_s\); by
\Cref{lem:smooth-pullback-intermediate-extension} these pullbacks are again
intermediate extensions of the corresponding pullbacks of \(E\). The difference
of the two cocycle morphisms restricts to zero on
\(\mathcal K_s\times\mathcal K_s\times O\), hence vanishes by
\Cref{lem:boundary-vanishing-smooth-pullback}. Thus the weak equivariant
structure is well defined.

We now verify the ordinary infinitesimal strong compatibility. Let
\(\xi\in\mathfrak k_K\). By
\Cref{lem:differentiation-equivariant-structures}, differentiating the weak
equivariant structure gives an operator \(A_\xi\) on \(j_{O,!+}(E)\). The
fundamental vector field gives the operator \(D_\xi\) through the
\(\mathscr D^\dagger_{\mathfrak X,\mathbb Q}\)-module structure. On the open
orbit \(O\), the equality \(A_\xi=D_\xi\) is the infinitesimal compatibility of
the equivariant isocrystal. Indeed, if
\[
  \nabla:E\longrightarrow E\otimes\Omega^1_O
\]
is the integrable connection corresponding to the isocrystal, then the
left \(\mathscr D^\dagger\)-module structure on the pair \((O,\overline O)\)
lets a vector field \(w\) act as \(\nabla_w\). Hence
\[
  D_\xi|_O=\nabla_{v_\xi}.
\]
The equivariant isomorphism \(\beta_E:a_O^*E\simeq p_O^*E\) is horizontal for
the pulled-back connections. Differentiating this horizontal isomorphism at the
identity of \(\mathcal K_s\) gives
\[
  A_\xi|_O=\nabla_{v_\xi}.
\]
Therefore \(A_\xi|_O=D_\xi|_O\).

The difference \(\Delta_\xi:=A_\xi-D_\xi\) is
\(\mathscr D^\dagger_{\mathfrak X,\mathbb Q}\)-linear. To see this, let
\(P\) be a local section of \(\mathscr D^\dagger_{\mathfrak X,\mathbb Q}\). Pull
back the \(\mathscr D^\dagger\)-linearity of \(\alpha_{j_{O,!+}(E)}\) along the
first-order point \(e_\xi\). The operator \(P\) is transformed into
\(P+\varepsilon\xi_{\mathscr D}(P)\), while the induced automorphism of the
module is \(\operatorname{id}+\varepsilon A_\xi\). Comparing the coefficient of
\(\varepsilon\) in
\[
  (\operatorname{id}+\varepsilon A_\xi)(P n)
  =(P+\varepsilon\xi_{\mathscr D}(P))
  (\operatorname{id}+\varepsilon A_\xi)(n)
\]
for a local section \(n\) gives
\[
  [A_\xi,P]=\xi_{\mathscr D}(P).
\]
By \Cref{lem:infinitesimal-functoriality-operators}, the action of the
fundamental vector field gives the same infinitesimal derivation on Berthelot
operators:
\[
  [D_\xi,P]=\xi_{\mathscr D}(P).
\]
Consequently
\begin{equation}
\label{eq:strong-middle-extension-commutator}
  [\Delta_\xi,P]
  =[A_\xi-D_\xi,P]
  =0.
\end{equation}
Thus
\[
  \Delta_\xi\in
  \operatorname{End}_{\mathscr D^\dagger_{\mathfrak X,\mathbb Q}}
  \bigl(j_{O,!+}(E)\bigr).
\]
Since \(\Delta_\xi|_O=0\), \Cref{lem:boundary-vanishing-endomorphism-difference}
gives \(\Delta_\xi=0\). Hence the weak equivariant structure is strong in the
ordinary infinitesimal sense.

The final assertion is not an additional theorem: it is exactly the definition
of the admissible set \(\Adm^{\Phi,\fstr}_{\mathcal K}(O,\overline O)\). If
\([E]\) belongs to that set, then the middle extension just constructed admits
the filtered finite-level realization required in
\Cref{def:strong-equivariance} and therefore lies in
\(F\text{-}\Coh^{\fstr}_{\mathfrak K}\).
\end{proof}

\end{document}
